\documentclass{ctexart}
\usepackage{amsmath,amssymb,amsthm}
\usepackage{geometry}
\usepackage{xcolor}
\usepackage{enumitem}
\usepackage{booktabs}
\usepackage{stmaryrd}
\geometry{a4paper, margin=1in}

\begin{document}

下面给出一个符号统一、定义完整的版本。全文中将 \textbf{sort} 翻译为``类型''，不译为``排序''。

\section*{Locale 上的束语义}

设 $X$ 是一个 locale，记其开集框架为

\[
\mathcal O(X).
\]

其底元和顶元分别记为

\[
\bot_X:=0_{\mathcal O(X)},
\qquad
\top_X:=1_{\mathcal O(X)}.
\]

对 $U,V\in\mathcal O(X)$，关系 $V\le U$ 表示 $V$ 是 $U$ 的子开集；$U\wedge V$ 表示交，$\bigvee_iU_i$ 表示并。

把 $\mathcal O(X)$ 看作 site：一个族 $(U_i\le U)_{i\in I}$ 是覆盖，当且仅当

\[
\bigvee_{i\in I}U_i=U.
\]

记 $X$ 上的层范畴为

\[
\mathrm{Sh}(X).
\]

\section*{1. 多类型一阶语言}

一个多类型一阶语言 $\mathcal L$ 包含：

\begin{enumerate}[leftmargin=2em]
    \item 类型符号的一族 $\mathcal S$；
    \item 对每个 $A\in\mathcal S$，一族类型为 $A$ 的变量；
    \item 函数符号
    \[
    f:A_1\times\cdots\times A_n\to B;
    \]
    \item 关系符号
    \[
    R\hookrightarrow A_1\times\cdots\times A_n.
    \]
\end{enumerate}

当 $n=0$ 时，函数符号是常元，关系符号是命题常元。

公式由原子公式、逻辑联结词

\[
\top,\quad\bot,\quad\wedge,\quad\vee,\quad\Rightarrow
\]

以及量词

\[
\forall,\quad\exists
\]

构成。否定是派生符号：

\[
\neg\varphi:=\varphi\Rightarrow\bot.
\]

\section*{2. $\mathrm{Sh}(X)$ 中的 $\mathcal L$-结构}

记

\[
\operatorname{Mod}_{\mathcal L}(\mathrm{Sh}(X))
\]

为所有定义在 $\mathrm{Sh}(X)$ 中的 $\mathcal L$-结构所成的类。

一个这样的结构记为

\[
\mathbf M\in
\operatorname{Mod}_{\mathcal L}(\mathrm{Sh}(X)).
\]

注意，不能写 $\mathbf M\in\mathrm{Sh}(X)$，因为 $\mathbf M$ 本身不是一个层，而是由若干层、层态射和子对象组成的结构。

\subsection*{2.1 类型的解释}

每个类型 $A\in\mathcal S$ 被解释为一个层

\[
\llbracket A\rrbracket_{\mathbf M}=F_A
\in\operatorname{Ob}(\mathrm{Sh}(X)).
\]

\subsection*{2.2 函数符号的解释}

对于函数符号

\[
f:A_1\times\cdots\times A_n\to B,
\]

给定一个层态射

\[
\llbracket f\rrbracket_{\mathbf M}:
F_{A_1}\times\cdots\times F_{A_n}
\longrightarrow F_B.
\]

在每个开集 $U\in\mathcal O(X)$ 上，它给出函数

\[
\llbracket f\rrbracket_{\mathbf M,U}:
F_{A_1}(U)\times\cdots\times F_{A_n}(U)
\longrightarrow F_B(U),
\]

并且与限制映射相容：

\[
\llbracket f\rrbracket_{\mathbf M,U}(a_1,\ldots,a_n)|_V
=
\llbracket f\rrbracket_{\mathbf M,V}
(a_1|_V,\ldots,a_n|_V).
\]

\subsection*{2.3 关系符号的解释}

对于关系符号

\[
R\hookrightarrow A_1\times\cdots\times A_n,
\]

给定一个子对象

\[
\llbracket R\rrbracket_{\mathbf M}
\hookrightarrow
F_{A_1}\times\cdots\times F_{A_n}.
\]

将该子对象视为一个子层，仍记为 $\llbracket R\rrbracket_{\mathbf M}$。于是，对每个 $U$，有包含关系

\[
\llbracket R\rrbracket_{\mathbf M}(U)
\subseteq
F_{A_1}(U)\times\cdots\times F_{A_n}(U).
\]

\section*{3. 上下文、赋值与项的解释}

设变量上下文为

\[
\Gamma=(x_1:A_1,\ldots,x_n:A_n).
\]

在模型 $\mathbf M$ 中，定义上下文对应的层

\[
F_\Gamma
:=
F_{A_1}\times\cdots\times F_{A_n}.
\]

对开集 $U\in\mathcal O(X)$，一个 $\Gamma$-赋值是一个元素

\[
\mathbf a=(a_1,\ldots,a_n)\in F_\Gamma(U),
\]

即

\[
a_i\in F_{A_i}(U).
\]

若 $V\le U$，则记

\[
\mathbf a|_V
:=
(a_1|_V,\ldots,a_n|_V).
\]

\subsection*{3.1 项的解释}

若项 $t$ 在上下文 $\Gamma$ 中的类型为 $B$，则对每个 $U$ 有

\[
\llbracket t\rrbracket_{\mathbf M,U}:
F_\Gamma(U)\longrightarrow F_B(U).
\]

其定义递归如下。

变量：

\[
\llbracket x_i\rrbracket_{\mathbf M,U}(\mathbf a)=a_i.
\]

若

\[
t=f(t_1,\ldots,t_k),
\]

其中

\[
f:A_1\times\cdots\times A_k\to B,
\]

则

\[
\llbracket t\rrbracket_{\mathbf M,U}(\mathbf a)
=
\llbracket f\rrbracket_{\mathbf M,U}
\bigl(
\llbracket t_1\rrbracket_{\mathbf M,U}(\mathbf a),
\ldots,
\llbracket t_k\rrbracket_{\mathbf M,U}(\mathbf a)
\bigr).
\]

项解释满足限制相容性：

\[
\llbracket t\rrbracket_{\mathbf M,U}(\mathbf a)|_V
=
\llbracket t\rrbracket_{\mathbf M,V}(\mathbf a|_V).
\]

\section*{4. Kripke--Joyal 迫近关系}

对公式 $\varphi$、开集 $U\in\mathcal O(X)$ 和赋值 $\mathbf a\in F_\Gamma(U)$，定义

\[
U\models_{\mathbf M}\varphi[\mathbf a].
\]

省略模型下标时，写作

\[
U\models\varphi[\mathbf a].
\]

\subsection*{4.1 原子公式}

\subsubsection*{等式}

若 $t_1,t_2$ 都是类型 $B$ 的项，则

\[
U\models t_1=t_2[\mathbf a]
\]

当且仅当

\[
\llbracket t_1\rrbracket_{\mathbf M,U}(\mathbf a)
=
\llbracket t_2\rrbracket_{\mathbf M,U}(\mathbf a)
\]

作为集合 $F_B(U)$ 中的元素相等。

\subsubsection*{关系}

若 $R$ 是 $n$-元关系符号，则

\[
U\models R(t_1,\ldots,t_n)[\mathbf a]
\]

当且仅当

\[
\bigl(
\llbracket t_1\rrbracket_{\mathbf M,U}(\mathbf a),
\ldots,
\llbracket t_n\rrbracket_{\mathbf M,U}(\mathbf a)
\bigr)
\in
\llbracket R\rrbracket_{\mathbf M}(U).
\]

\subsection*{4.2 命题联结词}

\subsubsection*{真}

\[
U\models\top[\mathbf a]
\]

总是成立。

\subsubsection*{假}

\[
U\models\bot[\mathbf a]
\quad\Longleftrightarrow\quad
U=\bot_X.
\]

也就是说，$\bot$ 只在 frame 的底元，即空开集上成立。

\subsubsection*{合取}

\[
U\models(\varphi\wedge\psi)[\mathbf a]
\]

当且仅当

\[
U\models\varphi[\mathbf a]
\quad\text{且}\quad
U\models\psi[\mathbf a].
\]

\subsubsection*{析取}

\[
U\models(\varphi\vee\psi)[\mathbf a]
\]

当且仅当存在覆盖

\[
U=\bigvee_{i\in I}U_i
\]

使得对每个 $i\in I$，至少有一个成立：

\[
U_i\models\varphi[\mathbf a|_{U_i}]
\]

或

\[
U_i\models\psi[\mathbf a|_{U_i}].
\]

不同的 $U_i$ 上可以由不同的析取项成立。

\subsubsection*{蕴涵}

\[
U\models(\varphi\Rightarrow\psi)[\mathbf a]
\]

当且仅当对所有 $V\le U$，都有

\[
V\models\varphi[\mathbf a|_V]
\quad\Longrightarrow\quad
V\models\psi[\mathbf a|_V].
\]

\subsubsection*{否定}

由于

\[
\neg\varphi:=\varphi\Rightarrow\bot,
\]

有

\[
U\models\neg\varphi[\mathbf a]
\]

当且仅当对所有 $V\le U$，

\[
V\models\varphi[\mathbf a|_V]
\quad\Longrightarrow\quad
V=\bot_X.
\]

\subsection*{4.3 量词}

设 $y$ 是类型为 $B$ 的变量，且 $F_B=\llbracket B\rrbracket_{\mathbf M}$。

\subsubsection*{存在量词}

\[
U\models\exists y^B\,\varphi(\mathbf a,y)
\]

当且仅当存在覆盖

\[
U=\bigvee_{i\in I}U_i
\]

以及截面

\[
b_i\in F_B(U_i)
\]

使得对每个 $i\in I$，

\[
U_i\models
\varphi(\mathbf a|_{U_i},b_i).
\]

因此，存在量词表达的是局部存在，而不一定要求在整个 $U$ 上存在一个统一见证。

\subsubsection*{全称量词}

\[
U\models\forall y^B\,\varphi(\mathbf a,y)
\]

当且仅当对所有 $V\le U$ 以及所有

\[
b\in F_B(V),
\]

都有

\[
V\models\varphi(\mathbf a|_V,b).
\]

\section*{5. 迫近关系的基本性质}

\subsection*{5.1 单调性}

若 $V\le U$，且

\[
U\models\varphi[\mathbf a],
\]

则

\[
V\models\varphi[\mathbf a|_V].
\]

\subsection*{5.2 局部性}

若

\[
U=\bigvee_{i\in I}U_i
\]

且对每个 $i\in I$ 都有

\[
U_i\models\varphi[\mathbf a|_{U_i}],
\]

则

\[
U\models\varphi[\mathbf a].
\]

单调性与局部性保证了迫近关系与层的限制、分片和胶合相容。

\section*{6. 公式的 frame 真值}

对固定模型 $\mathbf M$、公式 $\varphi$ 和赋值 $\mathbf a\in F_\Gamma(U)$，可以定义其真值开集为

\[
\|\varphi[\mathbf a]\|
:=
\bigvee
\left\{
V\le U
\;\middle|\;
V\models\varphi[\mathbf a|_V]
\right\}.
\]

特别地，对于闭公式 $\sigma$，其真值是 $\mathcal O(X)$ 中的元素：

\[
\|\sigma\|
=
\bigvee
\left\{
U\in\mathcal O(X)
\;\middle|\;
U\models_{\mathbf M}\sigma
\right\}.
\]

并且

\[
U\models_{\mathbf M}\sigma
\quad\Longleftrightarrow\quad
U\le \|\sigma\|.
\]

因此，公式的真值自然形成 frame $\mathcal O(X)$ 中的元素，而不是通常意义下的二值真值。

\section*{7. 全局满足与逻辑有效性}

若 $\sigma$ 是闭公式，则定义

\[
\mathrm{Sh}(X),\mathbf M\models\sigma
\]

当且仅当

\[
\top_X\models_{\mathbf M}\sigma,
\]

其中 $\top_X=1_{\mathcal O(X)}$ 是 frame $\mathcal O(X)$ 的顶元。

对所有模型都成立时，写作

\[
\mathrm{Sh}(X)\models\sigma,
\]

其含义是

\[
\forall\mathbf M\in
\operatorname{Mod}_{\mathcal L}(\mathrm{Sh}(X)),
\qquad
\top_X\models_{\mathbf M}\sigma.
\]

这里的 $\top_X$ 是 $\mathcal O(X)$ 的顶元，不是 $\mathrm{Sh}(X)$ 中的终对象。虽然二者都常被记为 $1$，但在本文中将它们区分为：

\[
\top_X\in\mathcal O(X),
\qquad
1_{\mathrm{Sh}(X)}\in\mathrm{Sh}(X).
\]

\section*{8. 直觉主义逻辑的可靠性}

若

\[
\vdash_{\mathrm{NJ}}\varphi
\]

表示 $\varphi$ 是直觉主义一阶逻辑中的定理，则对任意 locale $X$ 和任意

\[
\mathbf M\in
\operatorname{Mod}_{\mathcal L}(\mathrm{Sh}(X)),
\]

都有

\[
\top_X\models_{\mathbf M}\varphi.
\]

更一般地，若

\[
\Gamma\vdash_{\mathrm{NJ}}\varphi,
\]

则对任意开集 $U\in\mathcal O(X)$、任意模型

\[
\mathbf M\in
\operatorname{Mod}_{\mathcal L}(\mathrm{Sh}(X)),
\]

以及任意赋值

\[
\mathbf a\in F_\Gamma(U),
\]

都有

\[
U\models_{\mathbf M}\varphi[\mathbf a].
\]

因此正确的可靠性定理是

\[
\boxed{
\Gamma\vdash_{\mathrm{NJ}}\varphi
\Longrightarrow
\forall X\,
\forall\mathbf M\in
\operatorname{Mod}_{\mathcal L}(\mathrm{Sh}(X))
\,\forall U\in\mathcal O(X)
\,\forall\mathbf a\in F_\Gamma(U),
\quad
U\models_{\mathbf M}\varphi[\mathbf a]
}
\]

这里应称为\textbf{可靠性}（soundness），而不是``直觉主义完备性''。完备性是另一个方向，需要额外的语义类和证明系统假设。

\end{document}